\documentclass[12pt]{article}
\usepackage[utf8]{inputenc}
\usepackage[T1]{fontenc}
\usepackage{hyperref}
\title{Purepecha Language Resources for K-12 Curriculum Development}
\author{Compiled by Claude Sonnet 5 (claude-sonnet-5) for Andr\'as Kornai}
\date{}
\begin{document}
\maketitle

\section{Introduction}
Purepecha (P'urh\'epecha; ISO 639-3 tsz) is spoken by about 55,000 people in Michoac\'an, Mexico. It is a language isolate or the sole member of the Tarascan family. The search returned mostly scholarly and historical works, with little on school materials.

\section{Primary Educational Resources}
\begin{itemize}
\item Monz\'on's article ``Grammar for P'urhepecha speakers: dialectal variation and standard'' describes a grammar for native speakers, with Purepecha examples and Spanish explanations. It notes that the grammar has never been formally taught in schools: \url{https://revistaschilenas.uchile.cl/handle/2250/216755}
\item A University of Guadalajara paper on Purhepecha instruction at the Universidad Intercultural Ind\'igena de Michoac\'an: \url{https://riudg.udg.mx/handle/20.500.12104/94656}
\end{itemize}

\section{Linguistic and Reference Materials}
\begin{itemize}
\item Alberto Medina P\'erez, \emph{Gu\'ia ling\"u\'istica P'urh\'epecha} (2006), listed in Glottolog: \url{https://glottolog.org/resource/reference/id/147735}
\item Cristina Monz\'on, a bibliography of Purepecha vocabularies. It lists the dictionaries of Vel\'asquez Gallardo (1978), Medina P\'erez (1996) and Lathrop (1973, reissued 2007): \url{https://colmich.repositorioinstitucional.mx/jspui/bitstream/1016/1130/1/Monz%c3%b3nCristina2016Cap%c3%adtulo.pdf}
\item Juan Bautista de Lagunas, \emph{Arte y dictionario} (colonial period), a grammar and dictionary: \url{https://dare.uva.nl/id/9ea900a3-9879-45f2-ad5d-dc0d180cc9a6}
\item Omniglot overview of Purepecha and its writing: \url{https://www.omniglot.com/writing/purepecha.htm}
\end{itemize}

\section{Religious and Community Resources}
None found in search results.

\section{Digital and Audio Resources}
None found in search results.

\section{Curriculum Development Context}
Policy in Mexico calls for indigenous-language teaching, and published grammars and dictionaries exist. Classroom materials, audio and current school programmes still need to be found, ideally through Michoac\'an institutions.

\end{document}